% Author-supplied 24 September 2026 manuscript; provenance and local review are recorded separately.
\section{Zero-Hessian rigidity on full polycylinders}
\label{sec:supplied-p213}

The following argument is adapted from the supplied manuscript
\cite{SuppliedP21320260924}.


Let $K$ be a complete algebraically closed nontrivially valued non-Archimedean field. Fix $n\geq1$ and, for each $i$, either an open disk $|t_i-a_i|<R_i$ or an open annulus $\log|t_i-a_i|\in I_i$, with $a_i\in K$, $R_i>0$, and $I_i\subset\mathbb R$ a nonempty open interval. Their product is an ordinary open analytic domain $W$. Write $J$ for the annular indices.

The multiplication question for CLD currents asks whether $\mathrm d'\mathrm d''[U]=\mathrm d'\mathrm d''[V]$ for continuous plurisubharmonic potentials forces $\mathrm d'\mathrm d''(UT)=\mathrm d'\mathrm d''(VT)$ for every positive current $T$ closed under both differentials. On the domains above, the difference $H=U-V$ has a stronger property: the zero-current equation makes it an ordinary smooth function of zero smooth Hessian.

\begin{theorem}\label{thm:supplied-p213-polycylinder-rigidity}
Let $H$ be continuous on $W$ and suppose $\mathrm d'\mathrm d''[H]=0$ as an ordinary CLD current. There are real constants $c,b_i$, $i\in J$, such that
\begin{equation}\label{eq:supplied-p213-affine-form}
 H=c+\sum_{i\in J}b_i\log|t_i-a_i|\quad\text{on }W.
\end{equation}
In particular $\mathrm d'\mathrm d''(HT)=0$ for every positive doubly closed CLD current $T$ on $W$ for which continuous scalar multiplication is defined. Thus the equal-curvature multiplication identity holds on $W$ for every current in the continuous-potential convention.
\end{theorem}

The problem inventory already records fixed-domain strict polycylinder pairs among its established positive cases. The theorem gives a direct rigidity proof in the algebraically closed setting. The full-product hypothesis is part of the statement: an arbitrary ordinary neighborhood in an analytic space need not admit such coordinates, so the theorem does not settle multiplication by currents on general ordinary opens or over arbitrary ground fields.

The proof first restricts the zero-current equation to coordinate curves cut out by analytic equations. On each such curve, tests supported on a disk skeleton force constancy in attached disk directions. Density of rational points then makes $H$ a function of the annular radii alone. Finally, full torus tests force that function to be affine.

\subsection{Restriction to coordinate curves}\label{sec:supplied-p213-restriction}

We use two CLD facts in the 25 July 2025 revision: integration currents of boundaryless analytic spaces are doubly closed \cite[Proposition~7.7.9]{CLD25}, and the analytic Poincar\'e--Lelong formula has sign $\mathrm d'\mathrm d''[\log|f|]=[\operatorname{div}(f)]$ \cite[Theorem~16.7.3]{CLD25}. The following form of restriction keeps the argument within ordinary currents.

\begin{lemma}\label{lem:supplied-p213-cartier-restriction}
Let $X$ be a pure-dimensional boundaryless analytic space and let $H$ be continuous with $\mathrm d'\mathrm d''[H]=0$ on an ordinary open $\Omega$. If $Y\subset\Omega$ is locally obtained by successive effective Cartier cuts, none identically zero on a component of the preceding cycle, then $\mathrm d'\mathrm d''(H[Y])=0$. For a closed analytic embedding of a resulting pure-dimensional coordinate curve, the restriction $H|_Y$ has zero intrinsic Hessian current.
\end{lemma}

\begin{proof}
Induct on the number of cuts. Suppose $S$ is the positive doubly closed integration current of the preceding cycle and $\mathrm d'\mathrm d''(HS)=0$. For its next defining function $f$, choose a smooth convex function $\rho$ that equals $0$ far to the left and $t$ far to the right, with $0\leq\rho(t)-\max(t,0)\leq1$. Put $w_m=\rho(\log|f|+m)-m$, extended by the constant $-m$ near $f=0$. It is ordinary smooth, with positive smooth Hessian; on the unit locus that Hessian is $\rho''(\log|f|+m)\mathrm d'\log|f|\wedge\mathrm d''\log|f|$. Set $S_m=\mathrm d'\mathrm d''w_m\wedge S$. The smooth closed-form product rule gives
\[
 \mathrm d'\mathrm d''(HS_m)=\mathrm d'\mathrm d''w_m\wedge\mathrm d'\mathrm d''(HS)=0.
\]
The functions $w_m$ converge to $\log|f|$ away from its zero set and satisfy $|w_m|\leq|\log|f||+1$. The local integrability in Poincar\'e--Lelong gives $S_m\to[\operatorname{div}_S(f)]$ as ordinary currents. Positivity of $S_m$ and its weak convergence give uniform coefficient-mass bounds on compact test supports. Uniformly approximating $H$ there by smooth scalars yields $HS_m\to H[\operatorname{div}_S(f)]$, and the asserted vanishing follows.

For the intrinsic assertion, a compactly supported smooth test on a closed analytic subspace locally extends to an ambient test: lift the finitely many analytic units in its tropical chart, shrink until the lifts remain units, and apply compactly supported smooth cutoffs. Therefore pushforward of currents under the closed embedding is injective and commutes with the differentials. The ambient identity just proved implies the intrinsic identity.
\end{proof}

For later use, every coordinate fiber obtained by setting $t_j$ equal to $K$-rational constants for $j\neq i$ is such a Cartier intersection. Every cut has multiplicity one and is nonzero on the preceding coordinate product. Hence \cref{lem:supplied-p213-cartier-restriction} applies to $H$ on each coordinate disk or annulus.

\subsection{One-dimensional rigidity and rational points}\label{sec:supplied-p213-curves}

For a disk centered at $a$ and $b$ in it, write $\xi(b,r)$ for the Gauss point of the closed disk $|t-b|\leq r$. The CLD integration normalization evaluates a radial second-derivative test along these points.

\begin{lemma}\label{lem:supplied-p213-curve-rigidity}
Let $h$ be continuous with zero Hessian current on a full open disk $B=\{|t-a|<R\}$ or a full open annulus $A=\{\log|t-a|\in I\}$. On $B$ the function $h$ is constant. On $A$ it has the form $\alpha+\beta\log|t-a|$ with real $\alpha,\beta$.
\end{lemma}

\begin{proof}
Take $b\in B(K)$ and a smooth compactly supported function $g$ of $s=\log|t-b|$, supported below $\log R$. Extending $g(s)$ by zero at $b$ gives an ordinary smooth compactly supported CLD test on $B$. The one-dimensional integration formula and $\mathrm d'\mathrm d''[h]=0$ yield
\[
 \int_{\mathbb R}h(\xi(b,\mathrm e^s))g''(s)\,\mathrm ds=0.
\]
Thus $s\mapsto h(\xi(b,\mathrm e^s))$ is affine. Its finite limit $h(b)$ as $s\to-\infty$ makes it constant. Given $b,c\in B(K)$, choose $|b-c|<r<R$. Then $\xi(b,r)=\xi(c,r)$, so $h(b)=h(c)$.

For the annulus, apply the same tests on each attached disk $|t-b|<|b-a|$ for $b\in A(K)$. They show that $h(b)$ equals its value at the Gauss point of annular radius $|b-a|$. The latter value is affine as a function of $\log|t-a|$, since tests pulled back from the annular radius detect its distributional second derivative. Density of $K$-points, proved in \cref{lem:supplied-p213-rational-density}, extends the identities from rational points to the entire disk or annulus.
\end{proof}

\begin{lemma}\label{lem:supplied-p213-rational-density}
$K$-rational points are dense in every product domain $W$ above. For radii in $|K^\times|$, each multivariable Gauss point lies in the closure of rational points with those coordinate absolute values.
\end{lemma}

\begin{proof}
The topology of an affine analytic torus or space is generated by finitely many strict inequalities in absolute values of analytic polynomials on affinoid charts. At any point of a nonempty ordinary open, narrow such inequalities and the coordinate-radius conditions using bounds in $|K^\times|$. The value group of an algebraically closed nontrivially valued field is dense in $\mathbb R_{>0}$, so this produces a nonempty strict affinoid domain inside the given open. Its maximal ideals have finite residue extensions \cite[Corollary~5.19]{JonssonHorawa20}; since $K$ is algebraically closed, it contains a $K$-point.

For the assertion at a Gauss point with radii in $|K^\times|$, scale the coordinates to radius one. Given finitely many polynomial inequalities true at that point, the terms of maximal norm in each relevant polynomial have a nonzero Laurent-polynomial reduction. Choose residue coordinates off the finitely many zero loci of these reductions; the residue field is infinite. Lifting them gives a rational point with the prescribed radii and the same needed polynomial norms. Every neighborhood of the Gauss point therefore meets the asserted set.
\end{proof}

\subsection{Full products}\label{sec:supplied-p213-products}

We first treat a full polyannulus, where every factor is annular. Let $x_i=\log|t_i-a_i|$ and write $D=\prod_i I_i$; let $\sigma:D\to W$ be the multivariable Gauss section.

\begin{proposition}\label{prop:supplied-p213-polyannulus}
If $H$ is continuous on this full polyannulus and has zero Hessian current, then $H=F\circ x$ for a real affine function $F$ on $D$.
\end{proposition}

\begin{proof}
Fix all coordinates except $i$ at rational values. By \cref{lem:supplied-p213-cartier-restriction}, the remaining annular restriction has zero intrinsic Hessian, so \cref{lem:supplied-p213-curve-rigidity} says it depends only on $|t_i-a_i|$. Changing coordinates one at a time shows that $H$ is constant on rational points with a fixed norm vector $r\in(\log|K^\times|)^n$. By \cref{lem:supplied-p213-rational-density}, that value is $H(\sigma(r))$. Continuity and density of rational points then give $H=F\circ x$ everywhere, with $F=H\circ\sigma$ continuous.

For a compactly supported smooth superform $\eta$ of bidegree $(n-1,n-1)$ on $D$, its tropical pullback is compactly supported on $W$, because the preimage of a compact coordinate box is a compact closed polyannulus. The defining torus integration normalization, also valid with continuous scalar coefficient by uniform approximation, gives
\[
 0=\langle\mathrm d'\mathrm d''[H],x^*\eta\rangle
   =\int_D F\,\mathrm d'\mathrm d''\eta.
\]
Taking tests with a single coefficient complementary to $\mathrm d'x_i\wedge\mathrm d''x_j$ shows $\partial_i\partial_jF=0$ as a distribution for every $i,j$. A continuous function with zero distributional Hessian on the connected convex box $D$ is affine, as claimed.
\end{proof}

\begin{proof}[of \Cref{thm:supplied-p213-polycylinder-rigidity}]
For a general product, restrict $H$ to coordinate curves and apply \cref{lem:supplied-p213-curve-rigidity}. On $W(K)$ it is independent of each disk coordinate and, for each annular coordinate, depends only on its absolute value. Fix rational points in the disk factors. At every annular norm vector in $(\log|K^\times|)^J$, the fixed-norm closure assertion in \cref{lem:supplied-p213-rational-density} identifies the common value on rational points with the value at the corresponding annular Gauss point and the chosen disk points. Continuity of the Gauss section and density of rational points then give $H=F(r)$ on all of $W$, where $r$ records only the annular logarithmic radii. If $J$ is empty, $H$ is already constant.

Otherwise choose, inside every disk factor, a nonempty open annulus centered at $a_i$. Their product with all original annular factors is a full polyannulus. By \cref{prop:supplied-p213-polyannulus}, $H$ is affine in its full vector of logarithmic radii there. Its independence of the disk radii makes the corresponding coefficients zero. Since the annular radii still range over their entire original box, $F$ has the form in \eqref{eq:supplied-p213-affine-form} on $W$.

The right side of \eqref{eq:supplied-p213-affine-form} is ordinary CLD smooth and has zero smooth Hessian: each $t_i-a_i$ is an analytic unit on its annular factor. For every doubly closed $T$, the smooth product rule therefore yields $\mathrm d'\mathrm d''(HT)=0$. Applied to $H=U-V$, this is the equal-curvature multiplication identity.
\end{proof}
